% ID: SEL-D0007-BA-E01
% ESTADO: cotejada. Segunda pasada visual completa y borrador compilado revisados.
% Fuente: Selectividad/matematicas_aplicadas/ORD_TITULAR_MATEMÁTICAS APLICADAS CCSS-JUNIO-IMPRESO.pdf
% SHA-256: c42a605501a277fe7194b964ad92ad66047be9a62dc254fa2219d756d3d71317
% Página: 1; bloque A; ejercicio original 1. Puntuación total impresa.

\textbf{(2.5 puntos)} Una empresa de recambios industriales produce dos tipos
de baterías, $A$ y $B$. Su producción semanal debe ser de al menos 10 baterías
en total y el número de baterías de tipo $B$ no puede superar en más de
10 unidades a las fabricadas de tipo $A$. Cada batería de tipo $A$ tiene unos
gastos de producción de 150 euros y cada batería de tipo $B$ de 100 euros,
disponiendo de un máximo de 6000 euros a la semana para el coste total de
producción.

Si la empresa vende todo lo que produce y cada batería de tipo $A$ genera un
beneficio de 130 euros y la de tipo $B$ de 140 euros,
\textquestiondown{}cuántas baterías de cada tipo tendrán que producir a la
semana para que el beneficio total sea máximo? \textquestiondown{}Cuál es ese
beneficio?
